\section{Stan wewnętrzny i obserwacja agenta}
\label{sec:stan-i-obserwacja}

Ponieważ Ultimate Texas Hold’em jest grą z niepełną informacją, odwzorowanie tej
właściwości wymagało od mnie rozdzielenia pełnego stanu środowiska od informacji
udostępnianych agentowi.

W implementacji zastosowałem dwa odrębne modele danych: \texttt{RoundState}
przedstawiający pełny stan wewnętrzny rozdania oraz \texttt{UTHObservation}
reprezentujący obserwację dostępną agentowi.

Agent nie powinien i nie otrzymuje odwołania do obiektu
\texttt{RoundState}. Każda obserwacja jest tworzona na kanwie pełnego stanu przez
osobną funkcję projekcji, która wyciąga ze stanu wyłącznie informacje dozwolone
z perspektywy gracza.

\subsection{Pełny stan rozdania}
\label{subsec:round-state}

Obiekt \texttt{RoundState} zawiera wszystkie dane na temat rozdania. Zaweira on:

\begin{itemize}
    \item aktualną fazę gry,
    \item dwie karty własne gracza,
    \item dwie ukryte karty krupiera,
    \item odkryte karty wspólne,
    \item karty spalone,
    \item mnożnik wykonanego zakładu Play,
    \item końcowy wynik rozdania,
    \item szczegółowe rozliczenie zakładów,
    \item rezultat showdownu.
\end{itemize}

Naturalnie nie wszystkie pola są aktywne jednocześnie. W stanach pośrednich wynik
rozdania, rozliczenie, rezultat showdownu oraz mnożnik zakładu Play pozostają
nieokreślone. Pola te otrzymują wartości dopiero po przejściu do stanu
terminalnego.

Dla przykładu liczba kart wspólnych i spalonych jest zależna od aktualnej fazy.
Ta zaleznosc przedstawiam w tabeli~\ref{tab:karty-w-fazach}.

%TODO: dla komorek w kolumnach karty wspolne, karty spalone i faze - niech beda wysrodkwoane w pionie i poziomie
\begin{table}[htbp]
    \centering
    \caption{Liczba kart przechowywanych w stanie w poszczególnych fazach}
    \label{tab:karty-w-fazach}
    \begin{tabular}{
        |m{0.24\textwidth}
        |>{\centering\arraybackslash}m{0.13\textwidth}
        |>{\centering\arraybackslash}m{0.13\textwidth}
        |m{0.32\textwidth}|
    }
        \hline
        \textbf{Faza}
        &
        \textbf{Karty wspólne}
        &
        \textbf{Karty spalone}
        &
        \textbf{Znaczenie}
        \\
        \hline

        \texttt{PREFLOP}
        &
        0
        &
        0
        &
        Rozdano wyłącznie karty własne gracza i krupiera.
        \\
        \hline

        \texttt{FLOP}
        &
        3
        &
        1
        &
        Spalono jedną kartę i odkryto flop.
        \\
        \hline

        \texttt{RIVER}
        &
        5
        &
        2
        &
        Odkryto komplet kart wspólnych.
        \\
        \hline

        \texttt{TERMINAL}
        &
        5
        &
        2
        &
        Rozdanie zostało zakończone foldem albo showdownem.
        \\
        \hline
    \end{tabular}
\end{table}

Obiekt stanu zaimplementowałem jako niemodyfikowalny, wymusza to utworzenie nowej instancji \texttt{RoundState} 
przy każdym przejściu środowiska (zamiast modyfikowania istniejącego obiektu). Zdecydowałem się na takie rozwiązanie, żeby 
ograniczyć ryzyko, że obiekt stanu może potencjalnie zostać niewłaściwie zmodyfikowany, ale też takie rozwiązanie ułatwia tworzenie deterministycznych testów.

Podczas tworzenia stanu sprawdzane są między innymi:
\begin{itemize}
    \item poprawny typ fazy,
    \item obecność dokładnie dwóch kart gracza i dwóch kart krupiera,
    \item liczba kart wspólnych właściwa dla danej fazy,
    \item liczba kart spalonych właściwa dla danej fazy,
    \item brak zduplikowanych kart w całym stanie,
    \item zgodność pól końcowych z fazą rozdania.
\end{itemize}

Sprawdzane jest również, aby stan
niekońcowy nie zawierał wyniku ani rozliczenia, stan zakończony foldem
nie zawierał mnożnika Play ani rezultatu showdownu, natomiast stan
zakończony showdownem powinien zawierać zarówno mnożnik wykonanego zakładu,
jak i szczegóły ocenionych rąk.

\subsection{Obserwacja agenta}
\label{subsec:uth-observation}

Obiekt \texttt{UTHObservation} zawiera wyłącznie te dane, które mogą zostać
wykorzystane przez agenta podczas podejmowania decyzji. Są to:

\begin{itemize}
    \item aktualna faza gry,
    \item dwie karty własne gracza,
    \item dotychczas odkryte karty wspólne,
    \item zbiór legalnych akcji,
    \item mnożnik zakładu Play w obserwacji terminalnej.
\end{itemize}

W ramach tego obiektu umieściłem zbiór legalnych akcji, które agent może podjąć w danej fazie, 
dzięki temu zdejmuję z agenta konieczność odtwarzania zasad gry na podstawie fazy, co
ułatwia implementację agentów bazowych.

Obserwacja \texttt{UTHObservation}, podobnie jak pełny stan, jest obiektem niemodyfikowalnym.
Sprawdzana jest liczba widocznych kart, brak duplikatów oraz zgodność zbioru
legalnych akcji z aktualną fazą gry.

\subsection{Projekcja stanu na obserwację}
\label{subsec:projekcja-stanu}

Obiekt obserwacji dla agenta jest tworzony przez funkcję projekcji.

\begin{equation}
    O_t = \phi(S_t),
\end{equation}

gdzie \(S_t\) jest pełnym stanem środowiska w chwili \(t\), natomiast
\(O_t\) oznacza obserwację udostępnioną agentowi.

Dla warunków mojego środowiska funkcja projekcji ma postać

\begin{equation}
    \phi(S_t) =
    \left(
        p_t,
        C_t^{player},
        C_t^{community},
        A_t^{legal},
        m_t
    \right),
\end{equation}

gdzie:

\begin{itemize}
    \item \(p_t\) oznacza aktualną fazę gry,
    \item \(C_t^{player}\) oznacza karty własne gracza,
    \item \(C_t^{community}\) oznacza odkryte karty wspólne,
    \item \(A_t^{legal}\) oznacza zbiór legalnych akcji,
    \item \(m_t\) oznacza mnożnik zakładu Play, jeżeli został już wykonany.
\end{itemize}

Karty krupiera i karty spalone należą do \(S_t\), natomiast kolejność
kart w talii jest przechowywana przez wewnętrzne źródło kart. Żadna
z tych informacji nie jest elementem \(O_t\), dokładny zakres pól obu
obiektów przedstawiłem już w podrozdziale~\ref{subsec:round-state}
i~\ref{subsec:uth-observation}.

Ważne jest, że wynik rozdania, rozliczenie i szczegóły showdownu są po zakończeniu epizodu
zwracane przez obiekt \texttt{StepResult}, a nie przez samą obserwację.
Z racji, że te dane pojawiają się dopiero po wykonaniu ostatniej decyzji,
nie mogą wpłynąć na wybór akcji w bieżącym rozdaniu.

\section{Maszyna stanów i przestrzeń akcji}
\label{sec:maszyna-stanow}

Przebieg pojedynczego rozdania Ultimate Texas Hold’em odwzorowałem jako
maszynę stanów. Przejście do następnej fazy jest wyznaczone przez bieżącą
fazę i akcję, z kolei zawartość stanu, tym samym obserwacji kolejnego stanu zależy od kart
zwróconych przez źródło. Bieżąca faza określa zakres informacji
widocznych dla agenta, zbiór legalnych akcji oraz operacje wykonywane przez
środowisko po otrzymaniu decyzji. 

Przebieg rozdania można przedstawić jako sekwencję faz:
\begin{equation}
    \mathcal{P} =
    \{
        \texttt{PREFLOP},
        \texttt{FLOP},
        \texttt{RIVER},
        \texttt{TERMINAL}
    \}.
\end{equation}

Faza \texttt{TERMINAL} jest stanem absorbującym, tzn. nie ma dla niej żadnych legalnych
akcji, a rozdanie wchodzi
w nią wyłącznie w wyniku spasowania albo wykonania zakładu Play zakończonego
showdownem.

Metoda \texttt{reset()}, wywoływana na początku każdego rozdania, rozdaje po dwie karty
graczowi i krupierowi w kolejności \(P_1, D_1, P_2, D_2\), gdzie \(P\) to gracz a \(D\) to krupier, tworzy stan
\texttt{PREFLOP} i zwraca pierwszą obserwację. Karty wspólne i spalone nie są
jeszcze dobierane, to dzieje się w odpowiedniej fazie. Każde poprawne wykonanie metody \texttt{step()} przesuwa
rozdanie do kolejnej fazy albo do stanu terminalnego. Agent nie może pozostać
w tej samej fazie po wykonaniu akcji.

Pełna przestrzeń akcji składa się z sześciu decyzji domenowych:

\begin{equation}
    \mathcal{A} =
    \{
        \texttt{CHECK},
        \texttt{BET\_4X},
        \texttt{BET\_3X},
        \texttt{BET\_2X},
        \texttt{BET\_1X},
        \texttt{FOLD}
    \},
\end{equation}

Maszyna stanów z podziałem na fazy i zbiory legalnych akcji wygląda następująco:

\begin{equation}
    \mathcal{A}_{legal}(p) =
    \begin{cases}
        \{\texttt{CHECK}, \texttt{BET\_3X}, \texttt{BET\_4X}\},
        & p = \texttt{PREFLOP}, \\[1mm]
        \{\texttt{CHECK}, \texttt{BET\_2X}\},
        & p = \texttt{FLOP}, \\[1mm]
        \{\texttt{BET\_1X}, \texttt{FOLD}\},
        & p = \texttt{RIVER}, \\[1mm]
        \varnothing,
        & p = \texttt{TERMINAL}.
    \end{cases}
\end{equation}

Akcja \texttt{CHECK} odkłada decyzję o zakładzie Play i przesuwa rozdanie do
następnej fazy, spalając kartę oraz odkrywając flop albo turn i river. Każda
akcja typu \texttt{BET} ustala mnożnik Play wynikający z fazy: \(4\times\)
lub \(3\times\) przed flopem, \(2\times\) po flopie i \(1\times\) na riverze.

Po zakładzie Play środowisko odkrywa pozostałe karty wspólne, przeprowadza
showdown i przechodzi do stanu terminalnego. Akcja \texttt{FOLD} kończy
rozdanie na riverze bez showdownu.

Struktura maszyny stanów gwarantuje, że zakład Play może zostać wykonany najwyżej
raz w trakcie rozdania, każda akcja typu \texttt{BET} prowadzi bezpośrednio do stanu
terminalnego, a akcja \texttt{CHECK} jedynie przesuwa decyzję do późniejszej fazy.

Dla zilustrowania przebiegu rozdania posłużę się schematem przedstawionym na
rysunku~\ref{fig:maszyna-stanow-uth}.

\begin{figure}[htbp]
    \centering

    \begin{tikzpicture}[
        node distance=26mm and 28mm,
        phase/.style={
            draw,
            rounded corners,
            minimum width=29mm,
            minimum height=13mm,
            align=center,
            font=\small
        },
        terminal/.style={
            draw,
            very thick,
            double,
            rounded corners,
            minimum width=31mm,
            minimum height=14mm,
            align=center,
            font=\small
        },
        transition/.style={
            -{Latex[length=2.5mm]},
            thick
        },
        every node/.style={
            font=\small
        }
    ]

        \node[phase] (preflop) {
            \texttt{PREFLOP}
        };

        \node[phase, right=of preflop] (flop) {
            \texttt{FLOP}
        };

        \node[phase, right=of flop] (river) {
            \texttt{RIVER}
        };

        \node[terminal, below=of flop] (terminal) {
            \texttt{TERMINAL}
        };

        \draw[transition]
            (preflop)
            --
            node[above] {\texttt{CHECK}}
            (flop);

        \draw[transition]
            (flop)
            --
            node[above] {\texttt{CHECK}}
            (river);

        \draw[transition]
            (preflop)
            to[bend right=20]
            node[below left, align=center] {
                \texttt{BET\_3X}\\
                \texttt{BET\_4X}
            }
            (terminal);

        \draw[transition]
            (flop)
            --
            node[right] {\texttt{BET\_2X}}
            (terminal);

        \draw[transition]
            (river)
            to[bend left=20]
            node[below right, align=center] {
                \texttt{BET\_1X}\\
                \texttt{FOLD}
            }
            (terminal);

    \end{tikzpicture}

    \caption{Maszyna stanów pojedynczego rozdania UTH}
    \label{fig:maszyna-stanow-uth}
\end{figure}